\documentclass[11pt,a4paper]{article}

% --- Geometria e Layout ---
\usepackage[margin=2.5cm]{geometry}
\usepackage{fancyhdr}
\usepackage{setspace}
\onehalfspacing
\setlength{\headheight}{14.5pt}

% --- Codificação e Idioma ---
\usepackage[utf8]{inputenc}
\usepackage[T1]{fontenc}
\usepackage[brazil]{babel}

% --- Matemática e Símbolos (AMS) ---
\usepackage{amsmath,amssymb,amsfonts,amsthm}
\usepackage{mathtools}

% --- Cores, Gráficos, Vetorização e Tabelas ---
\usepackage{graphicx}
\usepackage{xcolor}
\usepackage{tikz}
\usetikzlibrary{arrows.meta,calc}
\usepackage{booktabs}
\usepackage{enumitem}

% --- Definição de Cores (Questão 170 — ENEM 2025) ---
\definecolor{caixaFundo}{HTML}{E6E7E8}
\definecolor{aguaFrente}{HTML}{55B3DF}
\definecolor{aguaTopo}{HTML}{A4D4EC}
\definecolor{aguaJato}{HTML}{2B9FD9}
\definecolor{raloFundo}{HTML}{6C77B2}
\definecolor{raloBorda}{HTML}{3B4269}
\definecolor{boiaSubmersa}{HTML}{5C6BA8}
\definecolor{textoCor}{HTML}{231F20}

% --- Definição de Cores (Questão 176 — ENEM 2024) ---
\definecolor{vermelhoTendencia}{HTML}{DD5151}
\definecolor{gradeCinza}{HTML}{939598}

% --- Macros TikZ para a Caixa de Descarga (Questão 170 — ENEM 2025) ---
\newcommand{\desenharCaixaBase}[2]{%
    % #1 = altura frontal da água (y_frente)
    % #2 = altura traseira da água (y_tras)
    \fill[caixaFundo] (0,0) -- (5.0,0) -- (5.0,4.0) -- (4.65,4.45) -- (0.35,4.45) -- (0,4.0) -- cycle;
    \fill[aguaFrente] (0,0) -- (5.0,0) -- (5.0,#1) -- (0,#1) -- cycle;
    \fill[aguaFrente] (0,#1) -- (5.0,#1) -- (4.65,#2) -- (0.35,#2) -- cycle;
    \fill[aguaTopo] (0,#1) -- (5.0,#1) -- (4.65,#2) -- (0.35,#2) -- cycle;
    \filldraw[fill=raloFundo, draw=raloBorda, line width=0.7pt] (2.5,0.42) ellipse (0.40cm and 0.14cm);
}

\newcommand{\desenharArestasCaixa}{%
    \begin{scope}[dashed, line width=0.85pt, dash pattern=on 3pt off 2pt, draw=textoCor]
        \draw (0,0) -- (0.35,0.85);
        \draw (5.0,0) -- (4.65,0.85);
        \draw (0.35,0.85) -- (4.65,0.85);
        \draw (0.35,0.85) -- (0.35,4.0);
        \draw (4.65,0.85) -- (4.65,4.0);
    \end{scope}
    \begin{scope}[line width=0.9pt, draw=textoCor, line join=round]
        \draw (0,0) rectangle (5.0,4.0);
        \draw (0,4.0) -- (0.35,4.45) -- (4.65,4.45) -- (5.0,4.0);
        \draw (0.35,4.0) -- (0.35,4.45);
        \draw (4.65,4.0) -- (4.65,4.45);
    \end{scope}
}

\newcommand{\desenharCotasBase}{%
    \begin{scope}[draw=textoCor]
        \draw[dashed, line width=0.6pt, dash pattern=on 2.5pt off 2pt] (0,0) -- (-1.15,0);
        \draw[dashed, line width=0.6pt, dash pattern=on 2.5pt off 2pt] (0,0.85) -- (-0.80,0.85);
        \draw[<->, >=Stealth, line width=0.6pt] (-1.08,0.03) -- (-0.75,0.82)
            node[midway, left=4pt, font=\sffamily\small, text=textoCor] {1,5 dm};
    \end{scope}
    \begin{scope}[draw=textoCor]
        \draw[dashed, line width=0.6pt, dash pattern=on 2.5pt off 2pt] (0,0) -- (-0.18,-0.58);
        \draw[dashed, line width=0.6pt, dash pattern=on 2.5pt off 2pt] (5.0,0) -- (5.18,-0.58);
        \draw[<->, >=Stealth, line width=0.6pt] (-0.15,-0.52) -- (5.15,-0.52)
            node[midway, fill=white, inner sep=4pt, font=\sffamily\small, text=textoCor] {2,5 dm};
    \end{scope}
}

% --- Configuração de Cabeçalho e Rodapé ---
\pagestyle{fancy}
\fancyhf{}
\lhead{Revisão de Matemática — ENEM}
\chead{}
\rhead{\thepage}
\rfoot{Data: \today}
\renewcommand{\headrulewidth}{0.4pt}
\renewcommand{\footrulewidth}{0.4pt}

% --- Início do Documento ---
\begin{document}

% =========================================================================
% Bloco 1 — Matemática básica e grandezas - 65 min
% =========================================================================

\section*{Questão 136 --- ENEM 2025}

Uma pessoa pretende instalar um kit de gás natural veicular (GNV) em seu carro. Na loja que escolheu para realizar a compra e instalação desse kit, havia cinco modelos de cilindro para armazenamento do gás, cujas capacidades, em metro cúbico, eram, respectivamente: $10$, $14$, $17$, $21$ e $25$. O preço do cilindro é proporcional à sua capacidade. Esse carro rodará $30\text{ km}$ diariamente, $7\text{ dias}$ por semana, e o consumo do GNV é de $1\text{ m}^3$ a cada $13\text{ km}$ rodados. A pessoa escolherá o modelo de cilindro de menor preço e que garanta apenas um abastecimento semanal.

Nessas condições, qual será a capacidade, em metro cúbico, do cilindro escolhido por essa pessoa?

\begin{enumerate}[label=(\Alph*)]
    \item $10$
    \item $14$
    \item $17$
    \item $21$
    \item $25$
\end{enumerate}

\section*{Questão 147 --- ENEM PPL 2024}

Um engenheiro civil organizou duas equipes de pedreiros, I e II, para a construção de um muro bastante extenso. Todos os pedreiros, de ambas as equipes, apresentam o mesmo rendimento por hora trabalhada frente à construção planejada. A equipe I era composta por $3$ pedreiros, que construíram $36\text{ metros quadrados}$ de muro, em $2\text{ dias}$, trabalhando $6\text{ horas por dia}$. A equipe II era composta por $5$ pedreiros, que trabalharam $8\text{ horas diárias}$, durante $6\text{ dias}$.

Quantos metros quadrados a equipe II construiu a mais do que a equipe I?

\begin{enumerate}[label=(\Alph*)]
    \item $144$
    \item $180$
    \item $204$
    \item $240$
    \item $276$
\end{enumerate}

\section*{Questão 148 --- ENEM PPL 2024}

Na primeira edição de uma Feira do Livro, a organização contou com a participação de dez voluntários, de mesma habilidade e experiência, que montaram cinco estandes de livros em uma semana, trabalhando todos os dias.

Na segunda edição da Feira do Livro, a organização dispôs de duas semanas para montar os estandes. Considere que todos os novos voluntários escolhidos tenham a mesma experiência, habilidade e produtividade dos voluntários da primeira edição da feira e que trabalharão todos os dias nessas duas semanas.

Considere duas situações para a segunda edição:
\begin{itemize}
    \item[\textbf{I:}] o número de voluntários permanecerá o mesmo;
    \item[\textbf{II:}] o número de estandes a serem montados permanecerá o mesmo.
\end{itemize}

Observe que o aumento no tempo de organização para montar os estandes, junto a uma das duas situações citadas, pode fazer o valor das outras variáveis (número de voluntários e de estandes) aumentar, diminuir ou permanecer igual.

O número de estandes na situação I e o número de voluntários na situação II, na segunda edição da feira, comparados aos números da primeira edição, serão, respectivamente,

\begin{enumerate}[label=(\Alph*)]
    \item maior e maior.
    \item maior e menor.
    \item igual e igual.
    \item igual e menor.
    \item menor e menor.
\end{enumerate}

\section*{Questão 149 --- ENEM 2023}

Um controlador de voo dispõe de um instrumento que descreve a altitude de uma aeronave em voo, em função da distância em solo. Essa distância em solo é a medida na horizontal entre o ponto de origem do voo até o ponto que representa a projeção ortogonal da posição da aeronave, em voo, no solo. Essas duas grandezas são dadas numa mesma unidade de medida.

A tela do instrumento representa proporcionalmente as dimensões reais das distâncias associadas ao voo. A figura apresenta a tela do instrumento depois de concluída a viagem de um avião, sendo a medida do lado de cada quadradinho da malha igual a $1\text{ cm}$.

\begin{center}
    \begin{tikzpicture}[scale=0.55, transform shape]
        % Borda da tela do instrumento
        \draw[thick] (-1.0, -1.2) rectangle (15.0, 4.8);

        % Malha quadriculada com linhas cinzas claras
        \draw[step=1.0, lightgray, thin] (0,0) grid (14,4);

        % Eixos coordenados com setas e rótulos
        \draw[-{Stealth}, thick] (0,0) -- (14.6,0) node[anchor=north west] {Distância};
        \draw[-{Stealth}, thick] (0,0) -- (0,4.6) node[anchor=south east] {Altitude};

        % Ponto e rótulo de Origem do voo
        \fill (0,0) circle (3.5pt);
        \node[align=center] at (0, -0.6) {Origem\\do voo};

        % Ponto e rótulo de Destino
        \fill (13,0) circle (3.5pt);
        \node[anchor=north] at (11.8, -0.2) {Destino};

        % Trajetória vetorial exata do voo
        \draw[very thick] (0,0) -- (2,1) -- (7,1) -- (10,0.3) -- (11,0.4) -- (13,0);
    \end{tikzpicture}
\end{center}

Essa tela apresenta os dados de um voo cuja maior altitude alcançada foi de $5\text{ km}$.

A escala em que essa tela representa as medidas reais é

\begin{enumerate}[label=(\Alph*)]
    \item $1 : 5$.
    \item $1 : 11$.
    \item $1 : 55$.
    \item $1 : 5\,000$.
    \item $1 : 500\,000$.
\end{enumerate}

\section*{Questão 140 --- ENEM 2024}

João e Felipe participaram, na escola, de uma maratona de matemática na qual, durante uma semana, resolveram $200$ questões cada. Nessa maratona, a porcentagem $P$ de acertos de cada participante é convertida em um conceito:

\begin{itemize}
    \item \textbf{insatisfatório:} se $0 \le P < 50$;
    \item \textbf{regular:} se $50 \le P < 60$;
    \item \textbf{bom:} se $60 \le P < 75$;
    \item \textbf{muito bom:} se $75 \le P < 90$;
    \item \textbf{excelente:} se $90 \le P \le 100$.
\end{itemize}

João acertou $75\%$ das questões da maratona e Felipe acertou $30\%$ a menos que a quantidade de questões que João acertou.

Os conceitos de João e Felipe foram, respectivamente,

\begin{enumerate}[label=(\Alph*)]
    \item muito bom e bom.
    \item muito bom e regular.
    \item muito bom e insatisfatório.
    \item bom e regular.
    \item bom e insatisfatório.
\end{enumerate}

\section*{Questão 178 --- ENEM 2015 (Caderno Azul)}

Para economizar em suas contas mensais de água, uma família de $10$ pessoas deseja construir um reservatório para armazenar a água captada das chuvas, que tenha capacidade suficiente para abastecer a família por $20$ dias. Cada pessoa da família consome, diariamente, $0{,}08\text{ m}^3$ de água.

Para que os objetivos da família sejam atingidos, a capacidade mínima, em litros, do reservatório a ser construído deve ser:

\begin{enumerate}[label=(\Alph*)]
    \item $16$
    \item $800$
    \item $1\,600$
    \item $8\,000$
    \item $16\,000$
\end{enumerate}

% =========================================================================
% Bloco 2 — Estatística e gráficos - 25 minutos
% =========================================================================

\section*{Questão 154 --- ENEM 2024}

Contratos de vários serviços disponíveis na internet apresentam uma quantidade excessiva de informações. Isso faz com que o tempo necessário para a leitura desses contratos possa ser longo.

O quadro apresenta uma amostra do tempo considerado necessário para a leitura completa do contrato de alguns serviços digitais.

\begin{center}
    \begin{tabular}{cc}
        \toprule
        \textbf{Serviço} & \textbf{Tempo (minuto)} \\
        \midrule
        A                & $36$                    \\
        B                & $17$                    \\
        C                & $27$                    \\
        D                & $13$                    \\
        E                & $13$                    \\
        F                & $13$                    \\
        \bottomrule
    \end{tabular}
\end{center}

\noindent {\footnotesize ROMERO, L. Não li e concordo. \textit{Superinteressante}, n. 307, ago. 2012 (adaptado).}

\vspace{0.4cm}

\noindent O tempo médio, em minuto, necessário para a leitura completa de um contrato de serviço dentre os listados no quadro é, com uma casa decimal, aproximadamente,

\begin{enumerate}[label=(\Alph*)]
    \item $13{,}0$.
    \item $15{,}0$.
    \item $19{,}8$.
    \item $20{,}0$.
    \item $23{,}3$.
\end{enumerate}

\section*{Questão 178 --- ENEM 2024}

A umidade relativa do ar é um dos indicadores utilizados na meteorologia para fazer previsões sobre o clima. O quadro apresenta as médias mensais, em porcentagem, da umidade relativa do ar em um período de seis meses consecutivos em uma cidade.

\begin{center}
    \begin{tabular}{|l|c|c|c|c|c|c|}
        \hline
        \textbf{Meses}        & Maio & Junho & Julho & Agosto & Setembro & Outubro \\ \hline
        \textbf{Umidade (\%)} & $66$ & $64$  & $54$  & $46$   & $60$     & $64$    \\ \hline
    \end{tabular}
\end{center}

Nessa cidade, a mediana desses dados, em porcentagem, da umidade relativa do ar no período considerado foi:

\begin{enumerate}[label=(\Alph*)]
    \item $56$.
    \item $58$.
    \item $59$.
    \item $60$.
    \item $62$.
\end{enumerate}

\section*{Questão 140 --- ENEM 2023}

Uma empresa de transporte faz regularmente um levantamento do número de viagens realizadas durante o dia por todos os $160$ motoristas cadastrados em seu aplicativo. Em um certo dia, foi gerado um relatório, por meio de um gráfico de barras, no qual se relacionaram a quantidade de motoristas com a quantidade de viagens realizadas até aquele instante do dia.

A distribuição da frequência da quantidade de viagens por motorista é dada por:
\begin{itemize}
    \item $1$ viagem: $10$ motoristas
    \item $2$ viagens: $10$ motoristas
    \item $3$ viagens: $55$ motoristas
    \item $4$ viagens: $25$ motoristas
    \item $5$ viagens: $0$ motoristas
    \item $6$ viagens: $50$ motoristas
    \item $7$ viagens: $10$ motoristas
\end{itemize}

Comparando os valores da média, da mediana e da moda da distribuição das quantidades de viagens realizadas pelos motoristas cadastrados nessa empresa, obtém-se:

\begin{enumerate}[label=(\Alph*)]
    \item $\text{mediana} = \text{média} < \text{moda}$
    \item $\text{mediana} = \text{moda} < \text{média}$
    \item $\text{mediana} < \text{média} < \text{moda}$
    \item $\text{moda} < \text{média} < \text{mediana}$
    \item $\text{moda} < \text{mediana} < \text{média}$
\end{enumerate}

% =========================================================================
% Bloco 3 — Geometria espacial e plana - 35 minutos
% =========================================================================

\section*{Questão 145 --- ENEM 2024}

O estádio do Maracanã passou por algumas modificações estruturais para a realização da Copa do Mundo de 2014, como, por exemplo, as dimensões do campo retangular. Para se adaptar aos padrões da Fifa, as dimensões do campo foram reduzidas de $110\text{ m} \times 75\text{ m}$ para $105\text{ m} \times 68\text{ m}$.

Em quantos metros quadrados a área do campo do Maracanã foi reduzida?

\begin{enumerate}[label=(\Alph*)]
    \item $24$
    \item $35$
    \item $555$
    \item $1\,110$
    \item $1\,145$
\end{enumerate}

\section*{Questão 142 --- ENEM 2023}

A água utilizada pelos $75$ moradores de um vilarejo provém de um reservatório de formato cilíndrico circular reto cujo raio da base mede $5\text{ metros}$, sempre abastecido no primeiro dia de cada mês por caminhões-pipa. Cada morador desse vilarejo consome, em média, $200\text{ litros}$ de água por dia.

No mês de junho de um determinado ano, o vilarejo festejou o dia do seu padroeiro e houve um gasto extra de água nos primeiros $20\text{ dias}$. Passado esse período, as pessoas verificaram a quantidade de água presente no reservatório e constataram que o nível da coluna de água estava em $1{,}5\text{ metro}$. Decidiram, então, fazer um racionamento de água durante os $10\text{ dias}$ seguintes. Considere $3$ como aproximação para $\pi$.

Qual é a quantidade mínima de água, em litro, que cada morador, em média, deverá economizar por dia, de modo que o reservatório não fique sem água nos próximos $10\text{ dias}$?

\begin{enumerate}[label=(\Alph*)]
    \item $50$
    \item $60$
    \item $80$
    \item $140$
    \item $150$
\end{enumerate}

\section*{Questão 170 --- ENEM 2025}

Uma caixa de descarga, acoplada a um vaso sanitário, tem a forma de paralelepípedo reto retângulo cujas dimensões internas da base são $2{,}5\text{ dm}$ e $1{,}5\text{ dm}$. Nessa caixa há uma boia que interrompe o abastecimento quando a altura da coluna de água atinge $2\text{ dm}$, conforme a figura.

\begin{center}
    \begin{tikzpicture}[scale=1.0]

        % =========================================================
        % FIGURA SUPERIOR: Caixa enchendo (boia abaixada)
        % =========================================================
        \begin{scope}[yshift=6.0cm]
            \desenharCaixaBase{0.62}{1.42}

            % Ondulações e jato de água
            \draw[aguaJato!80, line width=0.6pt] (4.48,1.02) ellipse (0.24cm and 0.08cm);
            \draw[aguaJato!90, line width=0.6pt] (4.48,1.02) ellipse (0.13cm and 0.045cm);
            \draw[aguaJato, line width=1.8pt, line cap=round] (4.48,3.45) -- (4.48,1.02);

            % Mecanismo da boia (abaixada)
            \draw[textoCor, line width=2.2pt] (4.45,3.45) -- (4.86,3.45);
            \draw[textoCor, line width=1.0pt] (4.46,3.45) -- (4.16,1.98);
            \fill[textoCor] (4.16,1.98) circle (0.22cm);

            \desenharArestasCaixa
            \desenharCotasBase
        \end{scope}

        % =========================================================
        % FIGURA INFERIOR: Abastecimento interrompido (altura 2 dm)
        % =========================================================
        \begin{scope}[yshift=0cm]
            \desenharCaixaBase{3.05}{3.62}

            % Mecanismo da boia (levantada na horizontal)
            \draw[textoCor, line width=2.2pt] (4.45,3.45) -- (4.86,3.45);
            \draw[textoCor, line width=1.0pt] (4.46,3.45) -- (2.90,3.45);

            % Esfera da boia (metade submersa e metade fora d'água)
            \begin{scope}
                \clip (2.90,3.45) circle (0.23cm);
                \fill[boiaSubmersa] (2.60,3.15) rectangle (3.20,3.40);
                \fill[textoCor] (2.60,3.40) rectangle (3.20,3.75);
                \draw[aguaTopo!70!boiaSubmersa, line width=0.6pt] (2.90,3.40) ellipse (0.23cm and 0.04cm);
            \end{scope}

            \desenharArestasCaixa
            \desenharCotasBase

            % Cota de 2 dm (altura da coluna de água à direita)
            \begin{scope}[draw=textoCor]
                \draw[dashed, line width=0.6pt, dash pattern=on 2.5pt off 2pt] (5.0,0) -- (6.05,0);
                \draw[dashed, line width=0.6pt, dash pattern=on 2.5pt off 2pt] (5.0,3.05) -- (6.05,3.05);
                \draw[<->, >=Stealth, line width=0.6pt] (5.92,0) -- (5.92,3.05)
                node[midway, fill=white, inner sep=4pt, font=\sffamily\small, text=textoCor] {2 dm};
            \end{scope}
        \end{scope}

    \end{tikzpicture}
\end{center}

A cada acionamento da descarga, todo o volume de água contida na caixa é despejado no vaso. Para reduzir o volume de água despejado a cada acionamento, uma pessoa colocará, no interior dessa caixa, garrafas de $300\text{ mL}$, cheias de areia e tampadas, de modo a ficarem submersas quando o abastecimento for interrompido.

Para garantir o funcionamento eficiente, o mínimo de água despejada a cada acionamento deve ser de $5\text{ L}$.

A quantidade máxima de garrafas que serão colocadas nessa caixa, garantindo um funcionamento eficiente, é igual a

\begin{enumerate}[label=(\Alph*)]
    \item $10$.
    \item $8$.
    \item $4$.
    \item $3$.
    \item $2$.
\end{enumerate}

% =========================================================================
% Bloco 4 — Funções e álgebra - 25 minutos
% =========================================================================

\section*{Questão 170 --- ENEM 2024}

Em uma região com grande incidência de terremotos, observou-se que dois terremotos ocorridos apresentaram magnitudes $M_1$ e $M_2$, medidos segundo a escala Richter, e liberaram energias iguais a $E_1$ e $E_2$, respectivamente. Entre os estudiosos do assunto, é conhecida uma expressão algébrica relacionando esses valores dada por
\[ M_2 - M_1 = \frac{2}{3} \log \left(\frac{E_2}{E_1}\right) \]

Estudos mais abrangentes observaram que o primeiro terremoto apresentou a magnitude $M_1 = 6{,}9$ e a energia liberada foi um décimo da observada no segundo terremoto.

O valor aproximado da magnitude $M_2$ do segundo terremoto, expresso com uma casa decimal, é igual a

\begin{enumerate}[label=(\Alph*)]
    \item $5{,}4$.
    \item $6{,}2$.
    \item $7{,}6$.
    \item $8{,}2$.
    \item $8{,}4$.
\end{enumerate}

\section*{Questão 176 --- ENEM 2024}

As receitas anuais obtidas por uma indústria no período de $2014$ a $2021$, em milhão de reais, foram registradas, por pontos, em um gráfico. Nele, também está representada a reta que descreve a tendência de evolução das receitas. Essa reta pode ser utilizada para estimar as receitas dos anos seguintes.

\begin{center}
    \begin{tikzpicture}[x=1.05cm, y=1.05cm]
        % Malha tracejada (horizontal e vertical)
        \begin{scope}[dashed, dash pattern=on 3pt off 3pt, draw=gradeCinza, line width=0.65pt]
            \foreach \y in {1,...,7} {
                    \draw (0,\y) -- (8.85,\y);
                }
            \foreach \x in {1,...,8} {
                    \draw (\x,0) -- (\x,7.70);
                }
        \end{scope}

        % Reta de tendência: passa exatamente por (2014, 4), (2017, 5) e (2020, 6)
        % Equação nas coordenadas da malha (x = ano - 2013): y = 11/3 + x/3
        \draw[vermelhoTendencia, line width=1.5pt] (0, {11/3}) -- (8.85, {11/3 + 8.85/3});

        % Pontos registrados no período de 2014 a 2021
        \foreach \x/\y in {
                1/3.45,
                2/5.00,
                3/4.87,
                4/4.56,
                5/5.59,
                6/6.00,
                7/5.17,
                8/6.71
            } {
                \fill[vermelhoTendencia] (\x,\y) circle (0.105);
            }

        % Eixos coordenados com setas triangulares sólidas
        \draw[-{Triangle[width=9.5pt, length=8.5pt]}, line width=1.15pt, draw=textoCor] (0,0) -- (8.85,0);
        \draw[-{Triangle[width=9.5pt, length=8.5pt]}, line width=1.15pt, draw=textoCor] (0,0) -- (0,7.75);

        % Rótulos numéricos do eixo vertical (Receita)
        \foreach \y in {0,...,7} {
                \node[anchor=east, font=\sffamily\small, text=textoCor] at (-0.12,\y) {\y};
            }

        % Rótulos numéricos do eixo horizontal (Ano)
        \foreach \x/\ano in {
                1/2014,
                2/2015,
                3/2016,
                4/2017,
                5/2018,
                6/2019,
                7/2020,
                8/2021
            } {
                \node[anchor=north, font=\sffamily\small, text=textoCor] at (\x,-0.12) {\ano};
            }

        % Títulos dos eixos
        \node[rotate=90, font=\bfseries, text=textoCor] at (-0.82, 3.85) {\textsf{Receita anual (em milhão de reais)}};
        \node[font=\bfseries, text=textoCor] at (4.5, -0.80) {\textsf{Ano}};
    \end{tikzpicture}
\end{center}

A estimativa da receita, em milhão de reais, dessa indústria, para o ano de $2026$, obtida a partir dessa reta de tendência, é

\begin{enumerate}[label=(\Alph*)]
    \item $7$.
    \item $8$.
    \item $9$.
    \item $10$.
    \item $11$.
\end{enumerate}

\section*{Questão 158 --- ENEM 2022}

Ao analisar os dados de uma epidemia em uma cidade, peritos obtiveram um modelo que avalia a quantidade de pessoas infectadas a cada mês, ao longo de um ano. O modelo é dado por $p(t) = -t^2 + 10t + 24$, sendo $t$ um número natural, variando de $1$ a $12$, que representa os meses do ano, e $p(t)$ a quantidade de pessoas infectadas no mês $t$ do ano. Para tentar diminuir o número de infectados no próximo ano, a Secretaria Municipal de Saúde decidiu intensificar a propaganda oficial sobre os cuidados com a epidemia. Foram apresentadas cinco propostas (I, II, III, IV e V), com diferentes períodos de intensificação das propagandas:

\begin{itemize}
    \item \textbf{I:} $1 \le t \le 2$;
    \item \textbf{II:} $3 \le t \le 4$;
    \item \textbf{III:} $5 \le t \le 6$;
    \item \textbf{IV:} $7 \le t \le 9$;
    \item \textbf{V:} $10 \le t \le 12$.
\end{itemize}

A sugestão dos peritos é que seja escolhida a proposta cujo período de intensificação da propaganda englobe o mês em que, segundo o modelo, há a maior quantidade de infectados. A sugestão foi aceita.

A proposta escolhida foi a

\begin{enumerate}[label=(\Alph*)]
    \item I.
    \item II.
    \item III.
    \item IV.
    \item V.
\end{enumerate}

\end{document}